\ifdefined\pdfsuppressptexinfo
  \pdfsuppressptexinfo=15
\fi
\documentclass[11pt]{article}

\usepackage[T1]{fontenc}
\usepackage[utf8]{inputenc}
\usepackage{lmodern}
\usepackage[margin=1in]{geometry}
\usepackage{amsmath,amssymb,amsthm,mathtools}
\usepackage{booktabs}
\usepackage{array}
\usepackage{enumitem}
\usepackage{microtype}
\usepackage{xcolor}
\usepackage{hyperref}

\hypersetup{
  colorlinks=true,
  linkcolor=blue!55!black,
  citecolor=blue!55!black,
  urlcolor=blue!55!black,
  pdftitle={Dirac-to-Tensor-Curl Eta Transfer on Round Spherical Space Forms with Affine-ADE Specializations},
  pdfauthor={Luqman Mahmood}
}

\newtheorem{theorem}{Theorem}[section]
\newtheorem{proposition}[theorem]{Proposition}
\newtheorem{lemma}[theorem]{Lemma}
\newtheorem{corollary}[theorem]{Corollary}
\theoremstyle{definition}
\newtheorem{definition}[theorem]{Definition}
\newtheorem{remark}[theorem]{Remark}
\newtheorem{example}[theorem]{Example}

\DeclareMathOperator{\Sym}{Sym}
\DeclareMathOperator{\Hom}{Hom}
\newcommand{\TT}{\mathrm{TT}}
\newcommand{\Z}{\mathbb Z}
\newcommand{\C}{\mathbb C}
\newcommand{\Irr}{\operatorname{Irr}}
\newcommand{\e}{\mathrm e}
\newcommand{\tr}{\operatorname{tr}}
\newcolumntype{L}[1]{>{\raggedright\arraybackslash}p{#1}}

\title{Dirac-to-Tensor-Curl Eta Transfer on Round Spherical Space Forms\\
with Affine-ADE Specializations}
\author{Luqman Mahmood\thanks{Copyright 2026 Luqman Mahmood. This paper is licensed under Creative Commons Attribution 4.0 International (CC BY 4.0); see \url{https://creativecommons.org/licenses/by/4.0/}. Software is excluded from this grant.}}
\date{1 October 2026}

\begin{document}

\maketitle

\begin{abstract}
Let $G<SO(4)$ be finite and effective, act freely on the unit round
three-sphere, and fix a lift of $G$ to
$\operatorname{Spin}(4)=SU(2)_L\times SU(2)_R$.  On the resulting spherical
space form, for every integer $J\geq1$ and every flat bundle associated to an
irreducible representation of $G$, we consider the eta invariant of the
helicity-signed transverse-traceless tensor-curl spectral family of rank $J$.
If $\boldsymbol\eta^D_C$ is the twisted spin-Dirac eta vector and
$A_m^L,A_m^R$ denote multiplication in the representation ring of $G$ by the
restricted symmetric powers of the two lifted defining representations, then
\[
 \boldsymbol\eta^{\TT}_J
 =-\frac12\left(A^L_{2J-1}+A^R_{2J-1}\right)
 \boldsymbol\eta^D_C.
\]
The proof combines a two-variable character identity with a continuation
lemma for finite-order Hurwitz--Lerch series and requires neither
commutativity of $G$ nor a one-sided action.  For one-sided subgroups of
$SU(2)$, the transfer specializes along the McKay correspondence: the eta
vectors satisfy a rank recurrence governed by the affine-ADE adjacency matrix
with an explicit linear-plus-periodic source.  We give closed formulas for
the cyclic and binary-dihedral families and specify the recurrence and source
periods for the exceptional binary polyhedral families.  The two-sided cyclic
case covers every lens space $L(p,q)$, while the nonabelian direct product
$Q_8\times\mathbb Z_3$ illustrates the general statement.  At rank one the
identity is consistent with the classical modulo-integers relation between
curl and Dirac eta invariants, with the Dirac sign specified below.  The contribution
here is the exact all-rank, flat-twisted transfer and its McKay-level
organization.
\end{abstract}

\section{Introduction}

Eta invariants measure the signed imbalance of a self-adjoint spectrum.
Their classical uses include index theory on manifolds with boundary,
rho-invariants, and phases of spectral determinants
\cite{APS1975,Donnelly1978,Gilkey1984}.  On spherical space forms, twisted
spin-Dirac eta invariants are well developed: character formulas,
Molien-series encodings, and complete spin-$\tfrac12$ tables for the finite
subgroups of $SU(2)$ are already in the literature
\cite{CisnerosMolina2001,Degeratu2009,Degeratu2001}.  The representation
theory behind the affine-ADE diagrams is likewise classical
\cite{McKay1980,Kostant1985,Rossmann2003}.

The operator studied here is different.  On the round sphere we use the
helicity-signed, projected transverse-traceless branch of the higher-spin
self-gradient family described by Homma and Homma--Tomihisa
\cite{Homma2001,HommaTomihisa2021}.  We denote its convention-fixed spectral
realization by $\mathcal C_J$.  At $J=1$ it is the classical coexact curl
operator; for $J\geq2$ the argument below uses the round-sphere module
decomposition and signed eigenvalues recorded in
\eqref{eq:tt-hilbert-space}--\eqref{eq:tt-sphere-branches}.  The rank-two
massive-gravity literature records the corresponding unsigned spectrum while
leaving its phase choice unresolved
\cite{Homma2001,HommaTomihisa2021,CastroLashkariMaloney2011,
BertinGrumillerVassilevichZojer2011}.  We retain the
helicity sign, analytically continue the imbalance, and organize all flat
twists by the representation ring of the deck group.

The rank-one member is the classical coexact curl sector, whose spectral and
eta-theoretic setting is already established \cite{APS1975,Bar2019}.  The
contribution begins at $J\geq2$ with the all-rank transfer and its McKay
organization; we do not claim novelty for curl itself.

The main result is not merely a list of exceptional-group values.  It is a
universal transfer from a classical spin-Dirac eta vector to every
integer-rank projected-TT eta vector on an oriented round spherical space
form with a chosen Spin(4) lift.  The proof is classwise and does not assume
that the deck group is cyclic or one-sided.  For one-sided finite subgroups of
$SU(2)$, this transfer turns the McKay graph into an eta-level computational
object and yields a sourced rank recurrence.  The cyclic, binary-dihedral,
and exceptional formulas then become explicit specializations of the general
identity.

The theorem concerns this round-sphere spectral family and its eta invariant;
it is not an operator intertwining or isospectrality statement.  The
affine-ADE labels organize the one-sided representation theory, without any
physical interpretation or selection of a group, representation, or
topology.

This paper is part of a series.  The lens-space conventions and closed-form
evaluations to which the cyclic corollary applies are developed in
\cite{MahmoodP1Collection2026}; the bivariate Chebyshev--McKay algebra
underlying the recurrences is treated in greater generality in
\cite{MahmoodP4Collection2026}; and the realization of $\mathcal C_J$ within a
four-dimensional boundary problem, together with the resulting APS
corollaries, appears in \cite{MahmoodP7Collection2026}.  The present proof is
self-contained and uses none of those companion results.

Section~\ref{sec:conventions} fixes the operator and continuation conventions.
Section~\ref{sec:transfer} proves the general transfer and its one-sided
specialization.  The remaining sections treat lens spaces, a mixed
nonabelian example, and the affine-ADE recurrences before discussing scope and
open problems.

\section{Operator, quotient, and twisting conventions}
\label{sec:conventions}

Identify $S^3$ with $SU(2)$.  Let $G<SO(4)$ be finite and effective, choose a
spin lift
\begin{equation}
 \tau=(\tau_L,\tau_R):G\longrightarrow SU(2)_L\times SU(2)_R,
 \qquad
 g\cdot x=\tau_L(g)x\tau_R(g)^{-1},
 \label{eq:spin4-action}
\end{equation}
and assume the action is free.  Equivalently, for every $g\ne e$,
\begin{equation}
 \chi_1(\tau_L(g))\ne\chi_1(\tau_R(g)).
 \label{eq:freeness-character}
\end{equation}
The ineffective central pair $(-1,-1)$ is therefore absent from $\tau(G)$.
Put $M_G=S^3/G$, with the unit round metric and orientation inherited
from the boundary of the standard oriented ball in $\mathbb C^2$.
The positive-left-heavy convention below agrees with
\cite{MahmoodP1Collection2026,MahmoodP7Collection2026}; in the latter's
full boundary operator the auxiliary finite character is subtracted.

\begin{proposition}[Lifts and spin structures]
\label{prop:spin-lifts}
Let $\pi:\operatorname{Spin}(4)\to SO(4)$ be the double covering.  A lift in
\eqref{eq:spin4-action} is a homomorphic section of the pulled-back central
extension
\[
 1\longrightarrow\{(1,1),(-1,-1)\}
 \longrightarrow\pi^{-1}(G)\longrightarrow G\longrightarrow1.
\]
For a free effective action, such lifts are equivalent to spin structures on
$M_G$.  They exist because every oriented three-manifold is spin, and the set
of lifts is a torsor for
\[
 H^1(M_G;\mathbb Z_2)\cong\Hom(G,\mathbb Z_2).
\]
Every lift is automatically injective.  If instead one starts with a
subgroup of $\operatorname{Spin}(4)$ containing $(-1,-1)$, its action on $S^3$ is
ineffective and the effective deck group is the quotient by that kernel.
\end{proposition}

\begin{proof}
A spin structure on the quotient pulls back to the unique spin structure on
$S^3$ and lifts the deck transformations compatibly with composition;
conversely, a homomorphic lift descends the spin bundle.  The ratio of two
lifts takes values in the central kernel and is a homomorphism
$G\to\mathbb Z_2$, while multiplying a lift by any such character produces
another lift.  Since $\pi\circ\tau$ is the inclusion of the effective group,
$\tau(g)=1$ implies $g=1$.  The last assertion is the kernel statement for
the $\operatorname{Spin}(4)$ action.  See also the spherical-space-form spin conventions in
\cite{Bar1992,Bar1996}.
\end{proof}

Write $V_n=\Sym^n\C^2$ for the irreducible $SU(2)$ representation of
highest weight $n$ and dimension $n+1$.  For integer $J\geq1$, put
\begin{equation}
 \mathcal H_J=\mathop{\widehat{\bigoplus}}_{n\geq0}
 \left(F^+_{J,n}\oplus F^-_{J,n}\right),
 \qquad
 F^+_{J,n}=V_{n+2J}\boxtimes V_n,
 \quad
 F^-_{J,n}=V_n\boxtimes V_{n+2J}.
 \label{eq:tt-hilbert-space}
\end{equation}
The completed sum uses the invariant Hilbert norms on its finite-dimensional
summands.  Define the signed spectral operator $\mathcal C_J$ by
\begin{align}
 \mathcal C_J|_{F^+_{J,n}}&=+(n+J+1),\nonumber\\
 \mathcal C_J|_{F^-_{J,n}}&=-(n+J+1),
 \qquad n\geq0.
 \label{eq:tt-sphere-branches}
\end{align}
Equations~\eqref{eq:tt-hilbert-space}--\eqref{eq:tt-sphere-branches} define the
representation-theoretic spectral realization used below.  The cited
higher-spin sources supply the round-sphere module decomposition and
eigenvalue normalization \cite{Homma2001,HommaTomihisa2021}.  For $J=1$ this
is the familiar coexact-curl branch.  For higher rank we refer to it as the
projected TT tensor-curl family in this convention-fixed sense; no extension
to a different metric or to an unprojected tensor operator is asserted.

\begin{proposition}[Analytic foundation]
\label{prop:analytic-foundation}
The operator $\mathcal C_J$ with domain
\[
 \operatorname{Dom}(\mathcal C_J)=
 \left\{v=(v_n^+,v_n^-):
 \sum_{n\geq0}(n+J+1)^2
 \bigl(\lVert v_n^+\rVert^2+\lVert v_n^-\rVert^2\bigr)<\infty\right\}
\]
is self-adjoint, has compact resolvent, and has no zero modes.  The left and
right $SU(2)$ actions preserve its domain and commute with it.  Hence the
induced operator on the multiplicity space
\begin{equation}
 \mathcal H_{J,\rho}
 =\Hom_G(V_\rho,\mathcal H_J)
 \cong(\mathcal H_J\otimes V_\rho^*)^G
 \label{eq:twisted-multiplicity-space}
\end{equation}
is self-adjoint with compact resolvent for every quotient satisfying
\eqref{eq:spin4-action}--\eqref{eq:freeness-character}.
\end{proposition}

\begin{proof}
The orthogonal summands in \eqref{eq:tt-hilbert-space} are finite-dimensional,
the displayed eigenvalues are real, and their absolute values tend to infinity.
The stated weighted-square-summability domain is therefore the maximal domain
of the diagonal self-adjoint operator, and its resolvent is compact.  Equivariance
is immediate because each summand is an $SU(2)\times SU(2)$ module on which
$\mathcal C_J$ is scalar.  Finite-group isotypic projections commute with the
operator and its resolvent.  Passing to the displayed multiplicity space,
rather than the full isotypic summand
$V_\rho\otimes\Hom_G(V_\rho,\mathcal H_J)$, is exactly the
copy-multiplicity normalization used below and introduces no extra factor of
$\dim\rho$.
\end{proof}

For $\rho\in\Irr(G)$, the associated-bundle normalization is the copy
multiplicity
\begin{equation}
 m_\rho(W)=\frac1{|G|}\sum_{g\in G}
 \overline{\chi_\rho(g)}\chi_W(g).
 \label{eq:copy-multiplicity}
\end{equation}
There is no additional factor of $\dim\rho$.  Replacing the associated-bundle
convention by its dual relabels $\rho$ by $\rho^*$ and changes no vector
identity when done consistently.

The full twisted eta invariant is
\begin{equation}
 \eta^{\TT}_{J,\rho}(s)=
 \sum_{n\geq0}
 \frac{m_\rho(F^+_{J,n})-m_\rho(F^-_{J,n})}
 {(n+J+1)^s},
 \qquad
 \eta^{\TT}_{J,\rho}=\eta^{\TT}_{J,\rho}(0).
 \label{eq:tt-eta-definition}
\end{equation}
The equality at $s=0$ means meromorphic continuation.  The operator has no
zero modes in these conventions.

The proof architecture is deliberately separated into four steps:
\begin{enumerate}[label=(\roman*)]
\item express each nonidentity class contribution as a finite
root-of-unity exponential--polynomial sum;
\item continue it by residue classes, proving that its value at $s=0$ is the
Abel value and is independent of the positive spectral shift;
\item compare the Dirac and TT rational generating functions at $t=1$; and
\item project the resulting class-function identity into irreducible
copy-multiplicity coordinates.
\end{enumerate}
This shift-independence step is what permits the Dirac origin
$n+\tfrac32$ and the TT origin $n+J+1$ to enter one eta identity.

\begin{lemma}[Finite-order Hurwitz--Lerch continuation]
\label{lem:hurwitz-continuation}
For a nonidentity class of a free effective finite action, each signed Dirac
or TT shell coefficient is a finite sum
\begin{equation}
 \sum_{\lambda}\sum_{k=0}^2 c_{\lambda,k}n^k\lambda^n,
 \qquad \lambda^{|G|}=1,\quad\lambda\ne1.
 \label{eq:finite-order-shell}
\end{equation}
Its eta continuation is a finite sum of shifted Hurwitz--Lerch functions
$\Phi(\lambda,s-j,a)$ and is regular at $s=0$.
Moreover,
\begin{equation}
 \Phi(\lambda,0,a)=\frac1{1-\lambda},
 \qquad\lambda\ne1,
 \label{eq:lerch-zero}
\end{equation}
independently of the positive spectral shift $a$.  Consequently the eta value
at zero equals the Abel value of the corresponding rational generating
function at $t=1$.
\end{lemma}

\begin{proof}
The two eigenvalues of each lifted factor are roots of unity.  A noncentral
$SU(2)$ character is a sum of two exponentials in $n$, while a Weyl limit at
$\pm1$ contributes at most one factor linear in $n$.  Products of the two
factor characters therefore have the form
\eqref{eq:finite-order-shell} with degree at most two.  Freeness gives
$z\ne w$ and $zw\ne1$ in eigenvalue coordinates for every nonidentity class;
equivalently, none of the four phases $zw,z/w,w/z,(zw)^{-1}$ that survives
the signed difference equals one.  If exactly one of $z,w$ is central, the
same conclusion follows by the corresponding single Weyl limit; freeness
keeps the remaining denominator nonzero.  The mixed fully central cases have
phase $-1$.
The identity shells cancel termwise, and the ineffective pair $(-1,-1)$ is
excluded.

If $q$ is the order of $\lambda$, decomposition into residue classes gives
\begin{equation}
 \Phi(\lambda,s,a)=q^{-s}\sum_{r=0}^{q-1}
 \lambda^r\zeta\!\left(s,\frac{a+r}{q}\right).
 \label{eq:lerch-hurwitz-decomposition}
\end{equation}
This formula, not the defining power series, is the meaning of
$\Phi(\lambda,s,a)$ at a root of unity on $|\lambda|=1$.
The possible pole at $s=1$ cancels because
$\sum_{r=0}^{q-1}\lambda^r=0$.  Using
$\zeta(0,x)=\tfrac12-x$ and
$\sum_{r=0}^{q-1}r\lambda^r=-q/(1-\lambda)$ gives
\eqref{eq:lerch-zero}.  For the polynomial weights use the finite identity
\[
 \sum_{n\geq0}n^k\lambda^n(n+a)^{-s}
 =\sum_{j=0}^k\binom{k}{j}(-a)^{k-j}\Phi(\lambda,s-j,a),
\]
initially where the series converges and then by continuation.  No
differentiation with respect to a root of unity is needed.  At nonpositive
integers the Bernoulli-polynomial formula for Hurwitz zeta gives the Abel
value of each polynomial-geometric series.  In particular, the value at
$s=0$ is the Abel value of $\sum n^k\lambda^n$, independent of $a$.
This proves regularity at zero and independence from the Dirac or TT shift.
\end{proof}

For the spin-Dirac control, use the branches
\begin{equation}
 D_n^+=V_n\boxtimes V_{n+1},\qquad
 D_n^-=V_{n+1}\boxtimes V_n,
 \label{eq:dirac-branches}
\end{equation}
with eigenvalues $\pm(n+3/2)$, and put
\begin{equation}
 \eta^D_{C,\rho}(s)=\sum_{n\geq0}
 \frac{m_\rho(D_n^+)-m_\rho(D_n^-)}{(n+3/2)^s}.
 \label{eq:dirac-eta-definition}
\end{equation}
On a one-sided quotient this is the sign printed in Cisneros--Molina's
spectral convention: the positive branch carries $V_n$ in $n+2$ copies and
the negative branch carries $V_{n+1}$ in $n+1$ copies
\cite{CisnerosMolina2001}.  We denote the resulting full twisted eta vector
by $\boldsymbol\eta^D_C$.  This fixes the Dirac sign convention used
throughout.

\section{General Spin(4) Dirac-to-tensor-curl transfer}
\label{sec:transfer}

For $g\in G$, write $z,z^{-1}$ and $w,w^{-1}$ for the defining eigenvalues
of $\tau_L(g)$ and $\tau_R(g)$.  For $\Re s>3$, define the equivariant
Dirichlet series
\begin{align}
 d_C(g;s)&=\sum_{n\geq0}
 \frac{\chi_n(z)\chi_{n+1}(w)-\chi_{n+1}(z)\chi_n(w)}
 {(n+3/2)^s},
 \label{eq:equivariant-dirac-series}\\
 t_J(g;s)&=\sum_{n\geq0}
 \frac{\chi_{n+2J}(z)\chi_n(w)-\chi_n(z)\chi_{n+2J}(w)}
 {(n+J+1)^s}.
 \label{eq:equivariant-tt-series}
\end{align}
Put $d_C(g)=d_C(g;0)$ and $\epsilon_J(g)=t_J(g;0)$ by the continuation of
Lemma~\ref{lem:hurwitz-continuation}.

\begin{lemma}[Two-variable class identity]
\label{lem:two-variable-identity}
For noncentral eigenvalue coordinates, set
\begin{equation}
 P_{a,b}(z,w;t)=\sum_{n\geq0}\chi_{n+a}(z)\chi_{n+b}(w)t^n.
 \label{eq:Pab-definition}
\end{equation}
For $|t|<1$,
\begin{align}
 P_{a,b}(z,w;t)=\frac1{(z-z^{-1})(w-w^{-1})}\bigg[&
 \frac{z^{a+1}w^{b+1}}{1-tzw}
 -\frac{z^{a+1}w^{-b-1}}{1-tz/w}\nonumber\\
 &-\frac{z^{-a-1}w^{b+1}}{1-tw/z}
 +\frac{z^{-a-1}w^{-b-1}}{1-t/(zw)}\bigg].
 \label{eq:Pab-rational}
\end{align}
At $t=1$, freeness makes the simplified denominators nonzero and
\begin{align}
 P_{0,1}-P_{1,0}
 &=-\frac{2zw}{(z-w)(zw-1)}
 =\frac{2}{\chi_1(w)-\chi_1(z)},
 \label{eq:dirac-rational-density}\\
 P_{2J,0}-P_{0,2J}
 &=\frac{zw\,[\chi_{2J-1}(z)+\chi_{2J-1}(w)]}
 {(z-w)(zw-1)}.
 \label{eq:tt-rational-density}
\end{align}
\end{lemma}

\begin{proof}
Insert
$\chi_n(z)=(z^{n+1}-z^{-n-1})/(z-z^{-1})$ into
\eqref{eq:Pab-definition} and sum the four geometric series, giving
\eqref{eq:Pab-rational}.  Subtraction and cancellation give
\eqref{eq:dirac-rational-density}--\eqref{eq:tt-rational-density}.  By
Lemma~\ref{lem:hurwitz-continuation}, these Abel values are the eta values at
zero despite the different positive spectral shifts.
\end{proof}

For $m\geq0$, let
\begin{equation}
 W_m^X=\Sym^m(\C^2)\circ\tau_X,
 \qquad
 (A_m^X)_{\rho\sigma}
 =\dim\Hom_G(\rho,W_m^X\otimes\sigma),
 \quad X=L,R.
 \label{eq:spin4-representation-matrices}
\end{equation}
Both matrices act on the representation ring of the full deck group $G$;
neither projection is required to be faithful.

\begin{theorem}[General free-Spin(4) Dirac-to-TT transfer]
\label{thm:spin4-transfer}
Under the hypotheses of Section~\ref{sec:conventions}, for every integer
$J\geq1$,
\begin{align}
 \epsilon_J(g)
 &=-\frac12\left[
 \chi_{2J-1}(\tau_L(g))+\chi_{2J-1}(\tau_R(g))
 \right]d_C(g),
 \label{eq:spin4-class-transfer}\\
 &\boxed{\boldsymbol\eta^{\TT}_J
 =-\frac12\left(A^L_{2J-1}+A^R_{2J-1}\right)
 \boldsymbol\eta^D_C.}
 \label{eq:spin4-transfer}
\end{align}
This is an identity of full eta invariants with copy-multiplicity
normalization.
\end{theorem}

Equation~\eqref{eq:spin4-transfer} relates meromorphically continued eta
invariants.  The Dirac and TT spectra lie on different arithmetic
progressions, so the formula is neither a shellwise relation nor an operator
intertwining or isospectrality statement.

\begin{proof}
For a noncentral pair,
\eqref{eq:spin4-class-transfer} is the quotient of
\eqref{eq:tt-rational-density} by
\eqref{eq:dirac-rational-density}.  Removable Weyl limits use
\begin{equation}
 \chi_m(1)=m+1,
 \qquad
\chi_m(-1)=(-1)^m(m+1).
\label{eq:weyl-central-values}
\end{equation}
If exactly one coordinate is central and the other is not, the corresponding
single Weyl limit exists; freeness keeps the simplified denominator nonzero,
so \eqref{eq:spin4-class-transfer} still holds.  The two fully central mixed
cases are
\begin{equation}
 (-1,1):\ d_C=\frac12,\ \epsilon_J=0,
 \qquad
 (1,-1):\ d_C=-\frac12,\ \epsilon_J=0.
 \label{eq:mixed-central-values}
\end{equation}
At $(1,1)$ the signed Dirac and TT shells vanish termwise.  The same
shellwise cancellation occurs at $(-1,-1)$, but that pair is excluded by
effectivity.

For clarity, all fully central possibilities are:
\begin{center}
\begin{tabular}{@{}ccccc@{}}
\toprule
$(\tau_L,\tau_R)$ & $d_C$ & $\epsilon_J$ &
$-\tfrac12(\chi_{2J-1}^L+\chi_{2J-1}^R)$ & status\\
\midrule
$(1,1)$ & $0$ & $0$ & $-2J$ & identity\\
$(-1,-1)$ & $0$ & $0$ & $+2J$ & ineffective, excluded\\
$(-1,1)$ & $\tfrac12$ & $0$ & $0$ & free mixed class allowed\\
$(1,-1)$ & $-\tfrac12$ & $0$ & $0$ & free mixed class allowed\\
\bottomrule
\end{tabular}
\end{center}

The irreducible characters are an orthonormal basis for class functions.
Multiplication by $\chi(W_m^X)$ has matrix $A_m^X$ in that basis, so projection
of \eqref{eq:spin4-class-transfer} using
\eqref{eq:copy-multiplicity} gives \eqref{eq:spin4-transfer}; no factor
$\dim\rho$ appears.
\end{proof}

\begin{remark}[Spin-lift covariance]
Two lifts differ by a sign character $\lambda:G\to\{\pm1\}$.  Simultaneous
central multiplication gives
\[
 d_C^{\lambda\tau}=\lambda d_C^\tau,
 \qquad
 W_m^{X,\lambda\tau}=\lambda^m\otimes W_m^{X,\tau}
 \quad(m\geq0),
\]
whereas every TT shell has even total central parity and is unchanged.  Thus
if $P_\lambda$ denotes the permutation of irreducible coordinates induced by
tensoring with $\lambda$, then
\[
 (\boldsymbol\eta^D_C)^{\lambda\tau}
 =P_\lambda(\boldsymbol\eta^D_C)^\tau,
 \qquad
 A_m^{X,\lambda\tau}=P_\lambda^m A_m^{X,\tau}
 =A_m^{X,\tau}P_\lambda^m,
 \qquad P_\lambda^2=I.
\]
Consequently, for general $m\geq0$,
\[
 A_m^{X,\lambda\tau}(\boldsymbol\eta^D_C)^{\lambda\tau}
 =P_\lambda^{m+1}A_m^{X,\tau}(\boldsymbol\eta^D_C)^\tau.
\]
For the odd powers $m=2J-1$ used in \eqref{eq:spin4-transfer}, the
factor $P_\lambda^{m+1}=P_\lambda^{2J}=I$ cancels.  This cancellation
is not asserted for even $m$.  In particular, $A_0^{X,\tau}=I$ for
either lift; replacing $P_\lambda^m$ by $P_\lambda$ would incorrectly
give $I=P_\lambda$ for a nontrivial sign character.
The integer-rank TT vector is lift-independent.
\end{remark}

\subsection{One-sided affine-ADE specialization}
\label{sec:one-sided-specialization}

Now let $G=\Gamma< SU(2)$ act from the left, so
$\tau_L(g)=g$ and $\tau_R(g)=1$.  Put
$M_\Gamma=\Gamma\backslash S^3$ and let
$\Delta=V_1$ be the defining representation.  The general eta definition
becomes
\begin{equation}
 \eta^{\TT}_{J,\rho}(s)=
 \sum_{n\geq0}
 \frac{(n+1)m_\rho(V_{n+2J})-(n+2J+1)m_\rho(V_n)}
 {(n+J+1)^s}.
 \label{eq:one-sided-tt-eta-definition}
\end{equation}
For a noncentral element with eigenvalues $u,u^{-1}$, set
\begin{equation}
 F_J(u)=u\left[
 \frac{u(u^{2J}-1)}{(1-u)^2}
 +\frac{u^{2J}-2J-1}{1-u}
 \right].
 \label{eq:FJ}
\end{equation}
Then
\begin{equation}
 \epsilon_J(g)=\frac{F_J(u)-F_J(u^{-1})}{u-u^{-1}}.
 \label{eq:epsilon-density}
\end{equation}
Character projection gives
\begin{equation}
 \eta^{\TT}_{J,\rho}(M_\Gamma)=\frac1{|\Gamma|}
 \sum_{g\in\Gamma}\overline{\chi_\rho(g)}\epsilon_J(g).
 \label{eq:class-transform}
\end{equation}
The identity and central $-1$ classes have value zero by signed shellwise
cancellation.

Define
\begin{equation}
 (A_m)_{\rho\sigma}
 =\dim\Hom_\Gamma(\rho,V_m\otimes\sigma).
 \label{eq:generalized-mckay}
\end{equation}
Thus $A_0=I$, $A_1$ is the ordinary McKay matrix, and
$A_{m+1}=A_1A_m-A_{m-1}$.  Since $A^R_{2J-1}=2JI$, the general theorem gives
the following corollary.

\begin{corollary}[One-sided Dirac-to-TT transfer]
\label{cor:one-sided-transfer}
For every finite $\Gamma< SU(2)$ and every integer $J\geq1$,
\begin{equation}
 \boxed{
 \boldsymbol\eta^{\TT}_J
 =-\left(JI+\frac12A_{2J-1}\right)
 \boldsymbol\eta^D_C.}
 \label{eq:transfer}
\end{equation}
Equivalently, with
\begin{equation}
 d_C(1)=0,
 \qquad
 d_C(g)=\frac{2}{2-\chi_\Delta(g)}\quad(g\ne1),
 \label{eq:dirac-density}
\end{equation}
one has
$\epsilon_J(g)=-(J+\chi_{2J-1}(g)/2)d_C(g)$, including the central values.
\end{corollary}

There is an equivalent finite-representation form that makes the identity
subtraction explicit.  Define
\begin{equation}
 P_J=\bigoplus_{r=0}^{2J-2}(2J-1-r)\Sym^r\Delta,
 \qquad
 \dim P_J=\frac{J(4J^2-1)}3.
 \label{eq:PJ}
\end{equation}

\begin{corollary}[Triangular representation form]
\label{cor:triangular}
For every $\rho\in\Irr(\Gamma)$,
\begin{equation}
 \boxed{
 \eta^{\TT}_{J,\rho}
 =-2J\eta^D_{C,\rho}
 +\langle\rho,P_J|_\Gamma\rangle
 -\frac{\dim\rho}{|\Gamma|}\frac{J(4J^2-1)}3.}
 \label{eq:triangular-transfer}
\end{equation}
\end{corollary}

\begin{proof}
The $SU(2)$ character recurrence gives
\begin{equation}
 (2-x)\chi_{P_J}=2J-\chi_{2J-1}.
 \label{eq:PJ-character}
\end{equation}
Indeed, for $a_r=2J-1-r$, the coefficient of $\chi_0$ in
$(2-x)\sum_ra_r\chi_r$ is $2a_0-a_1=2J$, every interior coefficient is
$2a_r-a_{r-1}-a_{r+1}=0$, and the terminal coefficient is $-1$.
Combining \eqref{eq:PJ-character} with \eqref{eq:dirac-density}, projecting,
and removing the identity contribution gives
\eqref{eq:triangular-transfer}.
\end{proof}

\begin{remark}[Sign conventions]
Some references use the sign-reversed Dirac operator $-D_C$, which reverses
the leading sign in the transfer, and reversing the two TT helicity branches
negates $\boldsymbol\eta^{\TT}_J$.  Every formula here uses the
Cisneros--Molina Dirac sign and the TT branch order fixed in
\eqref{eq:tt-sphere-branches}.
\end{remark}

\section{Lens-space specialization}
\label{sec:general-lens}

A general lens space has a mixed cyclic action and is not another affine-$A$
rung in the one-sided McKay classification.  It is, however, an immediate
cyclic specialization of Theorem~\ref{thm:spin4-transfer}.  The resulting
closed-form evaluations and detailed convention comparisons are developed in
\cite{MahmoodP1Collection2026}; here we record the transfer itself.

Write
\[
 L(p,q)=S^3/\langle\gamma\rangle,\qquad
 \gamma(z_1,z_2)=(\zeta z_1,\zeta^qz_2),\qquad
 \zeta=\e^{2\pi i/p},\quad \gcd(p,q)=1.
\]
Choose a lift bit $\ell$ satisfying
\begin{equation}
 q+1+\ell p\equiv0\pmod2
 \label{eq:lens-lift-condition}
\end{equation}
and put $\widetilde q=q+\ell p$.  The lifted left and right angles are
\begin{equation}
 \alpha=\frac{\pi(1+\widetilde q)}p,\qquad
 \beta=\frac{\pi(1-\widetilde q)}p.
 \label{eq:lens-lift-angles}
\end{equation}
Thus the two factors of a lifted deck class are represented by
$g_L=\e^{i\alpha}$ and $g_R=\e^{i\beta}$ in their defining tori.  Odd $p$ has
one valid lift; even $p$ has the two lifts $\widetilde q=q$ and $q+p$.

For this section the Dirac branches are
\[
 E_n^{D,+}=V_n\boxtimes V_{n+1},\qquad
 E_n^{D,-}=V_{n+1}\boxtimes V_n,
\]
and the TT branches are exactly
\eqref{eq:tt-sphere-branches}.  Let $d_C(g_L,g_R)$ denote the full
Cisneros-sign equivariant Dirac eta density for these branches.

\begin{corollary}[General lens-space transfer]
\label{thm:general-lens-transfer}
For every $p\geq2$, $\gcd(p,q)=1$, every valid spin lift, and every
$J\geq1$, a nonidentity deck class obeys
\begin{equation}
 \boxed{
 \epsilon_J(g_L,g_R)
 =-\frac12\left[
   \chi_{2J-1}(g_L)+\chi_{2J-1}(g_R)
 \right]d_C(g_L,g_R).}
 \label{eq:general-lens-class-transfer}
\end{equation}
After projection to all flat characters of $\Z_p$,
\begin{equation}
 \boxed{
 \boldsymbol\eta^{\TT}_J(L(p,q))
 =-\frac12\left(A^L_{2J-1}+A^R_{2J-1}\right)
 \boldsymbol\eta^D_C(L(p,q)).}
 \label{eq:general-lens-vector-transfer}
\end{equation}
The two matrices are the finite weight-convolution actions induced by the
left and right lifted circle homomorphisms.
\end{corollary}

\begin{proof}
Substitute the lifted torus elements into
\eqref{eq:spin4-class-transfer}.  Their characters are
\[
 \chi_{2J-1}(g_L)=\frac{\sin(2J\alpha)}{\sin\alpha},
 \qquad
 \chi_{2J-1}(g_R)=\frac{\sin(2J\beta)}{\sin\beta},
\]
with the Weyl limits
\eqref{eq:weyl-central-values} at central factors.  Fourier projection over
$\mathbb Z_p$ gives \eqref{eq:general-lens-vector-transfer}.
\end{proof}

The finite convolution can be written explicitly.  Set
\[
 s=\frac{1+\widetilde q}{2},\qquad
 t=\frac{1-\widetilde q}{2}\pmod p,
\]
and let $j$ run through the odd weights
$-(2J-1),-(2J-3),\ldots,2J-1$.  Then
\begin{equation}
 \eta^{\TT}_{J,a}
 =-\frac12\sum_j\left(
   \eta^D_{C,a-js}+\eta^D_{C,a-jt}
 \right).
 \label{eq:general-lens-convolution}
\end{equation}
For $q=1$ and $\ell=0$, one has $t=0$, so
$A^R_{2J-1}=2JI$ and the one-sided theorem
\eqref{eq:transfer} is recovered exactly.  Changing the even-order spin lift
shifts the Dirac vector and both odd kernels by $p/2$; these shifts cancel in
\eqref{eq:general-lens-convolution}, so the integer-rank TT vector is
lift-independent.

\section{A mixed nonabelian example}
\label{sec:mixed-q8c3}

A mixed nonabelian action is
\begin{equation}
 G=Q_8\times\mathbb Z_3,
 \qquad
 (a,u^r)\cdot x=a\,x\,u^{-r}.
 \label{eq:q8c3-action}
\end{equation}
The defining trace sets of the two projections are respectively
$\{2,-2,0\}$ and $\{2,-1\}$.  Their sole common trace is $2$, attained only
by the identity in each factor, so the action is free; $(-1,-1)$ is absent,
so it is effective.  The group has 15 complex
irreducible representations and four spin lifts.  Write
$A_0,A_1,A_2,A_3,D$ for the four one-dimensional and one two-dimensional
$Q_8$ irreps, and $c_0,c_1,c_2$ for the $\mathbb Z_3$ characters.  In the
manuscript convention, the rank-one values are
\begin{center}
\begin{tabular}{@{}lrrrr@{}}
\toprule
$Q_8$ irrep & $\eta^D_C(\rho\otimes c_0)$ & $\eta^D_C(\rho\otimes c_1)$
& $\eta^{\TT}_1(\rho\otimes c_0)$ & $\eta^{\TT}_1(\rho\otimes c_1)$ \\
\midrule
$A_0$ & $-89/144$ & $103/144$ & $-17/36$ & $-5/36$ \\
$A_1,A_2,A_3$ & $55/144$ & $-41/144$ & $19/36$ & $-5/36$ \\
$D$ & $-35/72$ & $13/72$ & $-4/9$ & $2/9$ \\
\bottomrule
\end{tabular}
\end{center}
The $c_2$ column equals the $c_1$ column.  These values illustrate the
theorem; they do not prove it.

\begin{example}[A non-self-dual coordinate and the matrix side]
\label{ex:nonselfdual-coordinate}
The characters $c_1$ and $c_2=c_1^*$ are genuinely complex and
non-self-dual.  The lifted defining representations are
\[
 W_1^L=D\otimes c_0,\qquad
 W_1^R=A_0\otimes(c_1\oplus c_2).
\]
With the row index denoting the target copy and the column index the source
copy in \eqref{eq:spin4-representation-matrices}, the
$A_0\otimes c_1$ component of the rank-one transfer is therefore
\[
 \eta^{\TT}_1(A_0c_1)
 =-\frac12\left[
 \eta^D_C(A_0c_0)+\eta^D_C(A_0c_2)+\eta^D_C(Dc_1)\right].
\]
The table gives
\[
 -\frac12\left[-\frac{89}{144}
 +\frac{103}{144}+\frac{13}{72}\right]
 =-\frac5{36},
\]
exactly the displayed TT coordinate.  This calculation fixes the
associated-bundle convention, shows why the multiplication matrix acts
without a transpose, and exhibits directly that no factor
$\dim(A_0c_1)$ or $\dim D$ is inserted.  Relabeling every coordinate by its
dual exchanges $c_1$ and $c_2$ and leaves the vector theorem unchanged.
\end{example}

\section{The affine-ADE source recurrence}
\label{sec:mckay}

Let $N=A_1$ be the McKay matrix and put
$E_J=(\eta^{\TT}_{J,\rho})_{\rho\in\Irr(\Gamma)}$.  For a noncentral class
with defining eigenvalues $u,u^{-1}$, define the recurrence residual
\begin{equation}
 r_J(g)=\epsilon_{J+1}(g)+\epsilon_{J-1}(g)
 -\chi_\Delta(g)\epsilon_J(g),
 \qquad J\geq2,
 \label{eq:source-residual-definition}
\end{equation}
and put $r_J(\pm1)=0$.  Substitution of
\eqref{eq:epsilon-density} and the $SU(2)$ character recurrence gives
\begin{equation}
 r_J(u)=-2J+
 \frac{(u^2+u+1)(u^{2J}-u^{-2J})}{u^2-1},
 \qquad u\neq\pm1,
 \label{eq:source-density}
\end{equation}
and set
\begin{equation}
 R_{J,\rho}=\frac1{|\Gamma|}\sum_{g\in\Gamma}
 \overline{\chi_\rho(g)}r_J(g).
 \label{eq:source-vector}
\end{equation}

\begin{theorem}[Sourced McKay recurrence]
\label{thm:mckay-source}
For $J\geq2$,
\begin{equation}
 \boxed{E_{J+1}+E_{J-1}-NE_J=R_J.}
 \label{eq:mckay-source}
\end{equation}
Moreover,
\begin{equation}
 \sum_{\rho}\dim(\rho)E_{J,\rho}=0,
 \qquad
 \sum_{\rho}\dim(\rho)R_{J,\rho}=0.
 \label{eq:regular-completeness}
\end{equation}
For every finite $\Gamma$, there is a vector $v$ and a periodic vector
$b_J$ such that
\begin{equation}
 R_J=Jv+b_J.
 \label{eq:linear-periodic}
\end{equation}
\end{theorem}

\begin{proof}
Project the defining identity
\eqref{eq:source-residual-definition} against the irreducible characters.
Multiplication by $\chi_\Delta$ becomes $N$, giving
\eqref{eq:mckay-source}; direct algebra in $u$ is the displayed
\eqref{eq:source-density}.  The regular character vanishes off the identity,
while $\epsilon_J$ and $r_J$ vanish at the identity, proving
\eqref{eq:regular-completeness}.  Every
noncentral eigenvalue $u$ is a root of unity.  The $u^{\pm2J}$ part is
therefore periodic in $J$, while the explicit first term is linear.  Finite
character transformation gives \eqref{eq:linear-periodic}.
\end{proof}

The shape ``linear plus periodic'' has a classical antecedent at the
multiplicity level for the spin-Dirac problem
\cite{CisnerosMolina2001}.  The claim here is the eta-vector lift for the
integer-rank projected-TT family, not novelty for the underlying
Clebsch--Gordan or McKay recurrence.

\section{Affine A: cyclic groups}
\label{sec:A}

For $\Gamma=\Z_p$ in the one-sided action, the irreducible representations
are the flat characters $\chi_c$, $c\in\Z/p\Z$, and the McKay graph is
$\widetilde A_{p-1}$.  Put $d=\min(c,p-c)$.  In the nonoverlap chamber
$p\geq4J-3$ for $J\geq2$ (and $p\geq2$ for $J=1$), the exact twisted formula
is
\begin{equation}
 \boxed{
 \eta^{\TT}_{J,c}(L(p,1))
 =-\frac{J}{3p}
 \left(p^2-6pd+6d^2+4J^2-2\right)
 +\left\lfloor\frac{\max(2J-d,0)^2}{4}\right\rfloor.}
 \label{eq:cyclic-stable}
\end{equation}
Outside this chamber the exact finite character formula
\eqref{eq:class-transform}, or equivalently the transfer theorem, remains
valid.  The first excluded row $p=4J-4$ differs only at the self-conjugate
character $c=p/2$, by $+1$.

\section{Affine D: every binary-dihedral group}
\label{sec:D}

Let
\begin{equation}
 \operatorname{Dic}_m=
 \langle a,x\mid a^{2m}=1,\ x^2=a^m,\ xax^{-1}=a^{-1}\rangle,
 \qquad m\geq2.
 \label{eq:dic-presentation}
\end{equation}
Its McKay graph is $\widetilde D_{m+2}$.  Write
\[
 h^{\mathrm{cyc}}_{J,c}
 =\eta^{\TT}_{J,c}(L(2m,1)),\qquad c\in\Z/(2m),
\]
for the cyclic value in the same orientation, generator, associated-bundle,
and full-eta convention.  When the value is represented through the Dirac
transfer, we use the canonical one-sided spin lift.

The four one-dimensional representations are
\[
 A_t(a)=1,\quad A_t(x)=t,\qquad t=\pm1,
\]
and
\[
 B_t(a)=-1,\quad B_t(x)=t,\qquad t=\pm i^m.
\]
Thus $t^2=(-1)^m$, as required by $x^2=a^m$.
The remaining irreducibles are two-dimensional:
\[
 V_r(a)=\begin{pmatrix}\zeta^r&0\\0&\zeta^{-r}\end{pmatrix},
 \qquad
 V_r(x)=\begin{pmatrix}0&(-1)^r\\1&0\end{pmatrix},
 \quad \zeta=\e^{\pi i/m},\quad1\leq r<m.
\]

\begin{theorem}[Binary-dihedral all-irrep corollary]
\label{thm:dic}
For every $m\geq2$ and $J\geq1$,
\begin{align}
 \eta^{\TT}_{J,A_t}&=\frac12h^{\mathrm{cyc}}_{J,0}-\frac{tJ}{2},
 \label{eq:dic-A}\\
 \eta^{\TT}_{J,B_t}&=\frac12h^{\mathrm{cyc}}_{J,m},
 \label{eq:dic-B}\\
 \eta^{\TT}_{J,V_r}&=h^{\mathrm{cyc}}_{J,r}.
 \label{eq:dic-V}
\end{align}
These formulas remain valid without a stable-chamber restriction.
\end{theorem}

\begin{proof}
The cyclic subgroup $\langle a\rangle$ has order $2m$.  The restrictions are
$A_t|_{\langle a\rangle}=\chi_0$,
$B_t|_{\langle a\rangle}=\chi_m$, and
$V_r|_{\langle a\rangle}=\chi_r\oplus\chi_{-r}$.
Every element of the coset $x\langle a\rangle$ has defining eigenvalues
$+i,-i$, where the TT density is $-J$.  Its character contribution is
$-tJ/2$ on $A_t$, cancels on $B_t$, and vanishes on $V_r$.
Finally $h^{\mathrm{cyc}}_{J,r}=h^{\mathrm{cyc}}_{J,2m-r}$, giving
\eqref{eq:dic-A}--\eqref{eq:dic-V}.
\end{proof}

At $m=2$, $\operatorname{Dic}_2=Q_8$ and the graph is
$\widetilde D_4$.  The two-dimensional defining-irrep eta vanishes for every
rank; the four line representations carry the nontrivial source.  For general
$m$, the source and rank shifts are explicit.  Put
\[
 Q_{J,c}=h^{\mathrm{cyc}}_{J+1,c}
 +h^{\mathrm{cyc}}_{J-1,c}
 -h^{\mathrm{cyc}}_{J,c-1}
 -h^{\mathrm{cyc}}_{J,c+1},
 \qquad c\pmod{2m}.
\]

\begin{corollary}[Binary-dihedral source]
\label{cor:dic-source}
Under the node labels of Theorem~\ref{thm:dic},
\begin{align}
 R_{J,A_t}&=\frac12Q_{J,0}-tJ,&
 R_{J,B_t}&=\frac12Q_{J,m},&
 R_{J,V_r}&=Q_{J,r}.
 \label{eq:dic-source}
\end{align}
The source after subtraction of its linear part has period dividing $m$.
\end{corollary}

\begin{corollary}[Binary-dihedral rank shift]
\label{cor:dic-rank-shift}
For the representatives $0\leq c<2m$,
\begin{equation}
 h^{\mathrm{cyc}}_{J+m,c}-h^{\mathrm{cyc}}_{J,c}
 =-\frac{4m^2-1}{12}+\frac{c(2m-c)}2+\frac{(-1)^c}{4}.
 \label{eq:cyclic-rank-shift}
\end{equation}
Consequently
\begin{align}
 \eta^{\TT}_{J+m,A_t}-\eta^{\TT}_{J,A_t}
   &=-\frac{m^2-1}{6}-\frac{tm}{2},\nonumber\\
 \eta^{\TT}_{J+m,B_t}-\eta^{\TT}_{J,B_t}
   &=\frac{2m^2+1+3(-1)^m}{24},\nonumber\\
 \eta^{\TT}_{J+m,V_r}-\eta^{\TT}_{J,V_r}
   &=-\frac{4m^2-1}{12}+\frac{r(2m-r)}2+\frac{(-1)^r}{4}.
 \label{eq:dic-rank-shift}
\end{align}
\end{corollary}

\begin{proof}
Write $\eta^D_c$ for the cyclic Dirac eta value in the canonical one-sided
lift.  The transfer identity~\eqref{eq:transfer} gives
\[
 h^{\mathrm{cyc}}_{J+m,c}-h^{\mathrm{cyc}}_{J,c}
 =-m\eta^D_c-\sum_{\substack{j\bmod 2m\\j\text{ odd}}}\eta^D_{c-j}.
\]
Indeed, the $2m$ added weights in the symmetric-power character cover each
odd residue modulo $2m$ twice.  For $0\leq c<2m$, character projection of
the one-sided Dirac density~\eqref{eq:dirac-density} evaluates to
\[
 \eta^D_c=\frac{4m^2-1}{12m}-\frac{c(2m-c)}{2m}.
\]
Summing this quadratic over the odd residues gives
$\sum_{j\text{ odd}}\eta^D_{c-j}=-(-1)^c/4$, where subscripts are reduced
modulo $2m$.  Substitution proves~\eqref{eq:cyclic-rank-shift};
Theorem~\ref{thm:dic} then gives the three binary-dihedral shifts.
\end{proof}

\section{Affine E: the exceptional binary polyhedral groups}
\label{sec:E}

For the binary tetrahedral, octahedral, and icosahedral groups, the same
class-transform and transfer theorems give the affine $E_6$, $E_7$, and $E_8$
eta vectors.  Table~\ref{tab:exceptional} records the explicit atlas scope.

\begin{table}[ht]
\centering
\small
\begin{tabular}{@{}lrrrL{0.34\textwidth}@{}}
\toprule
Group & Order & Irreps & Graph & Source structure \\
\midrule
$2T$ & 24 & 7 & $\widetilde E_6$ & linear plus vector period $3$ \\
$2O$ & 48 & 8 & $\widetilde E_7$ & linear plus vector period $12$ \\
$2I$ & 120 & 9 & $\widetilde E_8$ & linear plus vector period $15$ \\
\bottomrule
\end{tabular}
\caption{Exceptional affine-E families and the periods of their eta-vector
source terms.}
\label{tab:exceptional}
\end{table}

For $2I$, the conjugacy-class sizes are
\[
 (1,1,30,20,20,12,12,12,12),
\]
with defining traces
\[
 2,-2,0,1,-1,\varphi,\varphi^{-1},-\varphi^{-1},-\varphi.
\]
In the irrep order
\[
 (\mathrm{triv},\Delta,\Delta_2,\mathrm{3D},\mathrm{3D}_2,
  \mathrm{4S},\mathrm{4A5},\mathrm{5A5},\mathrm{6S}),
\]
the dimension vector $d=(1,2,2,3,3,4,4,5,6)$ satisfies $Nd=2d$.
The rank source has vector period $15$.  Analogous class formulas and periods
hold for $2T$ and $2O$.  Complete rational tables are omitted for space.  At
spin $\tfrac12$ such exceptional tables are classical
\cite{CisnerosMolina2001}; the present extension is the projected-TT
integer-rank eta tower.

\section{Numerical verification}

We checked the finite specializations with independent exact-arithmetic
implementations.  Direct spectral evaluation and the transfer formula agree
for every flat twist in forty lifted lens-space cases, for all 120 coordinates
of the $Q_8\times\mathbb Z_3$ example at ranks $1,\ldots,8$, and throughout the
cyclic, binary-dihedral, and binary-polyhedral tables used here.  These
computations check the sign, lift, and copy-multiplicity conventions; the
all-rank statements rest on the proofs above.

\section{Antecedents, novelty perimeter, and exclusions}
\label{sec:scope}

Flat-twisted eta invariants, equivariant fixed-point formulas, and the
rank-one coexact curl sector are classical
\cite{APS1975,Donnelly1978,Gilkey1984,Bar2019}.  So are spin-Dirac eta on
spherical space forms and its complete exceptional-group tables
\cite{Seade1981,Bar1992,Bar1996,Goette1997,CisnerosMolina2001}, the
Molien-series approach \cite{Degeratu2009,Degeratu2001}, and the
Clebsch--Gordan and McKay recurrences
\cite{McKay1980,Kostant1985,Rossmann2003}.  Higher-spin sphere spectra and
quotient-character methods provide the geometric and representation-theoretic
background \cite{Homma2001,HommaTomihisa2021,Dowker2007,
CastroLashkariMaloney2011}; charged spin-$\tfrac12$ and spin-$\tfrac32$
lens-space formulas provide related half-integer antecedents
\cite{BasileLeone2023,DebrayEtAl2023}.

The rank-one perimeter is especially strict.  Millson's equivariant curl
calculation is recorded here through B\"ar's full-text attribution, and
B\"ar treats curl spectra for general spherical space forms
\cite{Bar2019}.  Capoferri--Vassiliev state the classical APS-derived
relation, in their Dirac convention,
\[
 \eta_{\rm curl}(0)=-4\eta_{\rm Dirac}(0)+k,\qquad k\in\mathbb Z,
\]
on a general closed oriented three-manifold
\cite[Rem.~D.4(ii)--(iii)]{CapoferriVassiliev2023}.  Their Dirac operator
carries the opposite sign to the operator whose positive branch is
$V_n\boxtimes V_{n+1}$ in \eqref{eq:dirac-branches}.
In the present convention the same relation is
$\eta_{\rm curl}\equiv+4\eta^D_C\pmod{\mathbb Z}$.  For example, on
$L(p,1)$ with the canonical one-sided lift $\tau_R=1$, it reads
$\eta_{\rm curl}-4\eta^D_C=-(p-1)$.
The transfer identity is not a determination of the integer in the general
closed-manifold relation.  Thus the rank-one,
untwisted, modulo-integers shadow of
\eqref{eq:spin4-transfer} is not new.  Fixed-rank constructions of
higher-spin eta from twisted Dirac data in the Rarita--Schwinger anomaly
literature are also mechanism antecedents; the claim here is the exact,
mod-$\mathbb Z$-free, integer-rank vector identity with arbitrary flat
irreducible twists.

Eta invariants at flat connections also enter gauge Chern--Simons theory
classically \cite{Jeffrey1992}; we do not claim a new connection to that
theory.

Cisneros--Molina also gives a linear-plus-periodic multiplicity recurrence in
the spectral label.  Theorem~\ref{thm:mckay-source} must therefore be framed
as an eta-level lift to an operator family indexed by integer rank, not as the
first appearance of linear-periodic representation data.  Likewise, the new
claims concerning the exceptional families are confined to the projected-TT
higher-rank transfer and source recurrence at $J\geq2$.

To the best of our knowledge, the all-rank identities
\eqref{eq:spin4-transfer} and \eqref{eq:mckay-source} are new for $J\geq2$,
as is the genuinely mixed two-sided form of
Theorem~\ref{thm:spin4-transfer}.  This is not a claim of absolute priority.

The results are restricted to free, effective, orientation-preserving round
quotients with a specified spin lift.  They do not treat orbifolds, off-round
metrics, or a complete determinant or numerical APS index.
No physical interpretation is attached to these results.

\section{Discussion and outlook}

The transfer theorem changes the computational organization of the problem.
Without it, each rank and twist appears to require a new infinite signed
spectral continuation.  With it, the analytic work is concentrated in the
classical Dirac vector; every integer-rank TT eta vector follows by finite
representation-ring algebra.  The affine-ADE diagrams then organize the
calculation rather than merely label the deck groups.

The theorem supplies exact compact-manifold phase data and finite recurrences
that may enter a properly gauge-fixed determinant or topological field theory
calculation.  Constructing such a theory would additionally require ghosts,
zero modes, spectral cuts, and local counterterms, and lies beyond the present
spectral result.

The classwise identity now covers every free effective finite round action
with a chosen Spin(4) lift.  Three further directions are natural.  First, one
can turn the theorem into explicit irrep tables for the remaining mixed
nonabelian spherical-space-form families; this is an atlas problem, not a new
transfer theorem.  Second, one can study metric deformation away from the
round quotient.  Third, the eta transfer may be combined with determinant
magnitudes only after fixing ghosts, zero modes, spectral cuts, and local
counterterms.  Each requires separate analysis and none follows from the
present theorem.

A companion result realizes $\mathcal C_J$ naturally as the TT reducing
summand of the tangential operator of a four-dimensional chiral twisted-Dirac
problem on round spherical quotients.  That boundary realization is logically
independent of the eta-transfer proof and is not used here.  In particular,
it does not identify $\eta(\mathcal C_J)$ with the whole APS boundary
correction: the complementary boundary summand and its kernel remain separate
terms \cite{MahmoodP7Collection2026}.  The branchwise all-$s$ recurrences,
pole laws, heat expansions, and relative determinants are likewise developed
in the separate bivariate McKay manuscript rather than duplicated here
\cite{MahmoodP4Collection2026}.

\appendix

\section{Character identity behind the transfer}

For completeness, the triangular identity
\eqref{eq:PJ-character} can be checked without spectral analysis.  Write
\[
 \chi_{P_J}=\sum_{r=0}^{2J-2}(2J-1-r)\chi_r
\]
and use $x\chi_r=\chi_{r+1}+\chi_{r-1}$ with $\chi_{-1}=0$.  Every interior
coefficient cancels.  The two endpoints leave $2J\chi_0$ and
$-\chi_{2J-1}$.  This proof also shows that $P_J$ is finite and that
\[
 \dim P_J=\sum_{r=0}^{2J-2}(2J-1-r)(r+1)
 =\frac{J(4J^2-1)}3.
\]

\section{Regular completeness}

Let $\mathcal R=\bigoplus_\rho(\dim\rho)\rho$ be the regular representation.
Summing an eta vector with weights $\dim\rho$ is the same as taking the
regular character.  It vanishes off the identity; the signed sphere eta
density vanishes at the identity.  Hence
\[
 \sum_\rho\dim\rho\,\eta^{\TT}_{J,\rho}=0.
\]
The identical argument applies to the source vector.  This is the correct
completeness normalization; replacing copy multiplicity by isotypic dimension
would generally spoil it.

\section*{Acknowledgment of assistance}
AI-assisted research, computation, checking, and editorial support used
OpenAI Codex and Anthropic Claude. Luqman Mahmood directed the work
and is responsible for the claims and final text. These internal
checks are not external peer review.

\bibliographystyle{alpha}
\bibliography{references}

\end{document}
